% Main Content
\chapter{Introduction}

Local governance is fundamental to the upkeep of public infrastructure and the overall well-being of the community. To address public problems such as damaged roads, issues in waste management, street light failures, effective communication between citizens and local authorities is critical. Recently, the integration of e-governance systems and crowdsourcing has improve the citizen participation and streamline the reporting and resolution of public issues. 

However, most of the current systems rely on centralized architectures. This poses a challenge in dealing with data integrity, transparency, and user privacy. Citizens are reluctant to provide complaints about sensitive topics or criticism due to potential retaliation or the possibility of their personal info being abused. At the same time, the local government is faced with the challenge of verifying the authenticity of anonymous complaints raised by citizens.

Some of the key features of Blockchain technology are immutability, transparency, and decentralization. Those features can be used to address several limitations of conventional reporting systems. It is possible to design reporting platforms that protect user anonymity while ensuring the integrity and accountability of reported information by combining blockchain-based architectures and privacy-preserving mechanisms. This project examines the application of blockchain technology in developing a secure, community-driven reporting framework tailored to local governance contexts.


 % Local governance plays an important role in maintaining public infrastructure, ensuring community well-being, and addressing day-to-day issues faced by citizens. 
 
% Effective communication between citizens and local authorities is essential for identifying public problems such as damaged roads, waste management issues, streetlight failures, and other civic concerns. 


% In recent years, e-governance platforms and crowdsourced reporting systems have been introduced to improve citizen participation and streamline the reporting and resolution of public issues.


% However, despite these advancements, many existing reporting systems rely on centralized architectures that raise concerns regarding data integrity, transparency, and user privacy. 

% Citizens may be reluctant to report sensitive or critical issues due to fear of retaliation, misuse of personal data, or lack of trust in how reports are handled. At the same time, local governments face difficulties in verifying the authenticity and reliability of citizen-submitted reports, especially in open and anonymous reporting environments.
% Blockchain technology offers promising features such as immutability, transparency, and decentralized trust, which can address several limitations of conventional reporting systems. By leveraging blockchain-based architectures alongside privacy-preserving mechanisms, it is possible to design reporting platforms that protect user anonymity while ensuring the integrity and accountability of reported information. This project explores the application of blockchain technology to develop a secure, community-assisted reporting framework tailored for local governance contexts.



\section{Background}

% Crowdsourced reporting systems and e-governance platforms have been increasingly adopted worldwide in local governance. This has led to greater public participation and better service delivery. These tools let citizens report problems, give feedback, and join in decision making through digital platforms, which helps make government more inclusive and accountable. In developing countries like Sri Lanka, these systems can make governance more efficient by helping officials quickly spot and address public concerns.  

In the context of local governance, crowdsourced reporting systems and e-governance platforms are widely adopted. By using those systems, citizens can report issues, provide feedback and participate in decision making. It improves the inclusiveness and the accountability of every citizen. In developing countries like Sri Lanka, these systems have the potential to significantly improve the governance efficiency. Timely identification and the resolution of public conccerns are the key to governance efficiency. 

However, despite their advantages, the conventional e-governance and reporting systems are still highly centralized, making them prone to various threats such as data tampering, illegal access, and single points of failure. In addition, the ability to track behavior down to individual users could serve as a discouraging factor, particularly in reporting delicate matters, in a system that lacks anonymity. Blockchain-based systems have appeared as a feasible solution for managing shared data in a tamper-proof and transparent manner. The ability of permissioned blockchains to enable controlled participation while maintaining accountability through trusted validators makes them a desirable choice for governance applications. By using smart contracts, reporting workflows and opinion polling can be automated rather than relying on centralized intermediaries. 

While blockchain technology offers data integrity and transparency, it still poses challenges related to user privacy, content moderation, and trustworthiness among anonymous participants.  These challenges remain open research and implementation problems. Combining  blockchain with privacy-preserving identity mechanisms and \ac{AI}-assisted moderation can help address these limitations and supports the need for a comprehensive framework for secure and reliable community-assisted reporting.

% Crowdsourced reporting systems and e-governance platforms have been increasingly adopted worldwide to enhance public participation and improve service delivery in local governments. These systems allow citizens to report issues, provide feedback, and participate in decision making through digital platforms, thereby improving the inclusiveness and accountability. In developing countries such as Sri Lanka, these systems have the potential to significantly improve governance efficiency by enabling timely identification and resolution of public concerns.

% Despite their potential benefits, traditional e-governance and reporting platforms are predominantly centralized, making them vulnerable to data manipulation, unauthorized access, and single points of failure. Furthermore, the ability to trace back to individuals can discourage participation, particularly when reporting sensitive issues. The lack of anonymity also increases the risk of bias, suppression of reports, and reduced public trust.

% Blockchain-based systems have emerged as an alternative approach for managing shared data in a tamper-resistant and transparent manner. Permissioned blockchains, in particular, are well-suited for governance applications, as they allow controlled participation while maintaining accountability through trusted validators. Smart contracts enable the automation of reporting workflows, opinion polling, and issue resolution processes without reliance on centralized intermediaries. 

% However, while blockchain technology addresses data integrity and transparency, challenges related to user privacy, content moderation, and trustworthiness of anonymous participants remain open research and implementation problems. Integrating blockchain with privacy-preserving identity mechanisms and \ac{AI}-assisted moderation techniques offers a potential solution to these challenges, motivating the development of a comprehensive framework for secure and reliable community-assisted reporting.

\section{Problem Statement}

The existing frameworks of e-governance and crowdsourced reporting systems for local governance do not provide an integrated mechanism that can guarantee data integrity, user anonymity, and trustworthiness of citizen submitted reports. Centralized architectures are prone to unauthorized access, data tampering, and loss of transparency. At the same time, identity dependent systems discourage citizen engagement due to privacy concerns and fear of retaliation.

Although blockchain has been applied to the context of local government context, many blockchain-based reporting systems still fail to balance anonymity with mechanisms for assessing the reliability and quality of submitted reports. Furthermore, the lack of content moderation and trust evaluation increases the risk of misinformation, and malicious or low quality submissions. 

Therefore, there is a clear need for a privacy-preserving, blockchain-based reporting framework that enables anonymous citizen participation while ensuring data immutability, transparent workflows, and reliable evaluation of reported information. This project aims to address this gap by designing and prototyping a permissioned blockchain-based community-assisted reporting system that incorporates smart contracts, off-chain data storage, and \ac{AI}-assisted content moderation mechanism suitable for local governance environments.

% Existing e-governance and crowdsourced reporting systems for local governance lack an integrated mechanism to simultaneously ensure data integrity, user anonymity, and trustworthiness of citizen-submitted reports. Centralized architectures are prone to data tampering, unauthorized access, and loss of transparency. Also identity-dependent systems discourage citizen participation due to privacy concerns and fear of retaliation.




% Although blockchain technology has been explored in various governance applications, many existing blockchain-based reporting solutions either overlook user anonymity or fail to provide effective mechanisms for evaluating the reliability and quality of anonymous reports. Additionally, the absence of content moderation and trust evaluation mechanisms increases the risk of misinformation, malicious reporting, and low-quality submissions, reducing the practical usefulness of such systems for local authorities.

% Therefore, there is a clear need for a privacy-preserving, blockchain-based reporting framework that enables anonymous citizen participation while ensuring data immutability, transparent workflows, and reliable evaluation of reported information. This project aims to address this gap by designing and prototyping a permissioned blockchain-based community-assisted reporting system that incorporates smart contracts, off-chain data storage, and \ac{AI}-assisted content moderation mechanism suitable for local governance environments.

\section{Objectives and Scope}
\subsection{Aim}

The aim of this project is to design and prototype a blockchain-based, privacy-preserving, community-assisted reporting system for local governance that empowers citizens to report public issues and provide their opinions in a transparent, accountable, and secure manner.

\subsection{Objectives}

\begin{enumerate}
\item Design and implement a permissioned reporting blockchain using a \ac{PoA} consensus mechanism to ensure the integrity and immutability of citizen-submitted reports.

% \item Design and deploy smart contracts to automate public issue reporting workflows, including ticket creation, community voting, issue resolution, closure, and public opinion polling.

\item Design and deploy smart contracts to automate public issue reporting workflows, including ticket creation, issue resolution and closure, community voting as feedback on resolved issues, and separate opinion polling mechanisms for public consultation.

\item Develop user-friendly interfaces that enable citizens and authorized government officials to interact with the blockchain for reporting issues and tracking their resolution status.

\item Design and integrate an AI-based content moderation system to detect and filter inappropriate textual and visual content in user-submitted reports.

% \item Analyze user participation behavior and implement an \ac{AI}-assisted trustworthiness evaluation mechanism to improve the reliability of community feedback while preserving user anonymity.

\end{enumerate}


\subsection{Scope}

\begin{enumerate}
   \item A prototype of the proposed blockchain will be implemented with a minimum of three full nodes. The \ac{PoA} consensus mechanism will be used to enable efficient and controlled transaction validation within a permissioned network, while excluding deployment of public, permissionless blockchains or globally distributed ledger infrastructures. 
   
   \item User authentication and identity decoupling are managed through a simulated identity service, and the design and implementation of a fully decentralized Self-Sovereign Identity (SSI) blockchain are explicitly excluded from the scope of this project.
   
   % \item The functional scope of the blockchain layer is strictly limited to managing the lifecycle of reporting tickets  and facilitating public opinion polls, while other e-Governance services such as taxation, civil registration, or licensing systems are not considered. 

   % \item The blockchain is restricted to storing transaction metadata and cryptographic hashes, and the storage of high-capacity data such as images, videos, or detailed \ac{GPS} logs directly on-chain is excluded, with such data instead handled through off-chain storage mechanisms such as \ac{IPFS}.
   
   \item The integration of \ac{AI} within the system is limited to content moderation using existing machine learning models, libraries, or external Application Programming Interfaces (\ac{API}s), and the development of custom, from-scratch deep learning architectures is not within the scope of this work.
   
   % \item The trust evaluation mechanism assesses user reliability solely based on historical reporting behavior and peer-validation voting outcomes, while the incorporation of external data sources such as social media activity, biometric information, or financial credit records is excluded.     % is this needed?
   
   \item Transaction validation authority within the blockchain network is restricted to pre-approved institutional stakeholders, including local government entities and recognized non-governmental organizations, and does not permit anonymous public participation or Proof-of-Work (\ac{PoW}) based mining.
   

   \item The project emphasizes functional validation, system evaluation, and performance analysis within a controlled experimental environment, and does not address real-world legal compliance, regulatory adoption, legislative reform, or large-scale integration with existing national government information systems.

   \item The user interface development is limited to a responsive, web-based \ac{DApp} accessible via standard mobile and desktop web browsers, and explicitly excludes the development of native mobile applications (iOS or Android) or cross-platform native frameworks.

\end{enumerate}







